\section{Experimental Setup}

We design experiments to answer the following questions:

\textbf{RQ1:} Does fix-impact-ratio discriminate strategic debugging (targeting root causes) from sequential trial-and-error (addressing failures one-by-one)?

\textbf{RQ2:} Do agents cluster errors by shared characteristics (measured via clustering coefficient $> 0.3$) compared to random ordering?

\textbf{RQ3:} Does error clustering enable prioritization of high-impact fixes (modifications resolving $\geq 2$ test failures simultaneously)?

\textbf{RQ4:} Do agents transfer learned patterns to held-out test cases (passing unseen tests without error messages) at rates exceeding random baseline?

\subsection{Datasets}

We evaluate on synthetic multi-test programming problems to establish metric sensitivity before real-world deployment.

\textbf{Mock Synthetic Dataset:}
\begin{itemize}
    \item \textbf{Problems:} 50 programming tasks (sum, max, array operations, string manipulation, basic algorithms)
    \item \textbf{Test Cases:} 15--25 per problem (mean 20.3, median 19)
    \item \textbf{Error Types:} Balanced distribution---syntax (24.3\%), runtime (25.3\%), logic (27.4\%), edge\_case (23.0\%)
    \item \textbf{Rationale:} Controlled environment validates whether metrics \emph{can} detect strategic behavior when engineered. Balanced error types ensure clustering signal detectable. Mock validation precedes production deployment with GPT-4 + Codeforces.
\end{itemize}

\subsection{Baselines}

\textbf{Random Sampling:} Agent generates code variations without error feedback (temperature=0.7), no access to test failure information. Establishes null model for fix-impact-ratio (expected ratio $\approx 1.0$ if no strategic behavior).

\textbf{Sequential Trial-and-Error:} Agent addresses test failures one-by-one in test index order, no error clustering or root-cause analysis. Baseline for clustering coefficient (expected $\approx 1.0$ under random error ordering).

Baselines represent non-strategic approaches, enabling discrimination of clustering and prioritization behaviors.

\subsection{Implementation Details}

\textbf{Mock Agent Architecture:}
\begin{itemize}
    \item \textbf{Framework:} Python simulation with controlled parameters (\texttt{clustering\_strength}, \texttt{fix\_success\_rate}, \texttt{pattern\_memory\_enabled})
    \item \textbf{Clustering Simulation:} \texttt{clustering\_strength=0.5} yields 50\% probability of clustering same-type errors consecutively vs random ordering
    \item \textbf{Fix Success:} \texttt{fix\_success\_rate=0.6} (baseline) vs \texttt{0.7} (proposed method), simulating root-cause prioritization effectiveness
    \item \textbf{Iteration Budget:} \texttt{max\_iterations=10} per problem, \texttt{temperature=0.7}
    \item \textbf{Seed:} Fixed random seed (\texttt{seed=1}) for reproducibility
\end{itemize}

\textbf{Hyperparameters (by hypothesis):}
\begin{itemize}
    \item \textbf{h-e1:} $N=10$ problems, controlled trajectories (strategic vs sequential)
    \item \textbf{h-m1:} $N=50$ problems, \texttt{clustering\_strength=0.5}, permutation test with 1000 samples
    \item \textbf{h-m2:} $N=50$ problems, \texttt{baseline\_fix\_success=0.6}, \texttt{proposed\_fix\_success=0.7}, \texttt{cluster\_bonus\_probability=0.6}
    \item \textbf{h-m3:} $N=50$ problems, \texttt{revealed\_fraction=0.5}, \texttt{pattern\_memory\_enabled=True}, \texttt{max\_iterations=10}
\end{itemize}

\textbf{Compute Resources:} Local execution, $<1$ hour total runtime (mock agents, no GPU required)

\textbf{Reproducibility:} Code available in \texttt{h-\{id\}/code/} directories. Implementation uses controlled parameters to validate metric sensitivity, not production agent behavior.

\subsection{Evaluation Metrics}

\textbf{Fix-Impact-Ratio (FIR):} Tests passed per modification, $\text{FIR} = \frac{1}{k} \sum_{i=1}^{k} \Delta_{\text{passing}}(m_i)$. Strategic agents achieve ratio $> 2.0$ (one fix resolves 2+ failures), sequential agents $\approx 1.0$. Evaluated via t-test or Mann-Whitney U ($p < 0.05$). \textbf{Why:} Directly measures debugging efficiency---core metric validating strategic vs sequential discrimination (RQ1).

\textbf{Clustering Coefficient (CC):} Consecutive same-type errors vs expected under random permutation, $\text{CC} = \frac{C_{\text{observed}}}{C_{\text{random}}}$. Coefficient $> 0.3$ indicates clustering above chance. Evaluated via permutation test (1000 shuffles, $p < 0.05$). \textbf{Why:} Tests whether agents recognize error patterns (RQ2), validating mechanism Stage 1 (clustering enables prioritization).

\textbf{High-Impact Proportion:} Percentage of modifications with $\Delta_{\text{passing}} \geq 2$ (multi-test fixes). Proposed method expected $>$ baseline. Evaluated via proportion test ($p < 0.05$). \textbf{Why:} Tests prioritization effectiveness (RQ3), validating mechanism Stage 2 (clustering $\rightarrow$ high-impact fixes).

\textbf{Held-Out Test Slope Ratio:} Agent slope (held-out pass rate vs iteration) divided by random-mutation slope. Ratio $> 1.5$ with $p < 0.05$ indicates transfer learning. \textbf{Why:} Tests pattern transfer to unseen tests without error messages (RQ4), attempting to validate mechanism Stage 3 (clustering $\rightarrow$ transfer).

\textbf{Effect Size:} Cohen's d for fix-impact-ratio (0.2=small, 0.5=medium, 0.8=large). Large effect size ($d > 0.8$) demonstrates strong discriminative power.

Statistical significance threshold $p < 0.05$ for all tests. Permutation tests control for problem-specific error distributions (clustering coefficient, held-out slope). Directional tests used where appropriate (high-impact proportion: proposed $>$ baseline).
